%----------------------------------------------------------------
% Progress report for the supervisors (Tiziana Cattai, Pierluigi Locatelli)
% Methods + results: Part 1 (water networks) and Part 2 (power, traffic).
%----------------------------------------------------------------
\documentclass[11pt,a4paper]{article}

\usepackage[T1]{fontenc}
\usepackage[utf8]{inputenc}
\usepackage{mathptmx}
\usepackage[a4paper,top=2.4cm,bottom=2.4cm,left=2.6cm,right=2.6cm]{geometry}
\usepackage[english]{babel}
\usepackage{amsmath,amssymb}
\usepackage{graphicx}
\usepackage{booktabs}
\usepackage{tabularx}
\usepackage{array}
\usepackage{xcolor}
\usepackage{titlesec}
\usepackage[hidelinks]{hyperref}
\usepackage{enumitem}
\setlist{noitemsep,topsep=2pt}

\titlespacing*{\section}{0pt}{10pt}{4pt}
\titlespacing*{\subsection}{0pt}{7pt}{3pt}
\titleformat{\section}{\normalfont\large\bfseries}{\thesection}{0.6em}{}
\titleformat{\subsection}{\normalfont\bfseries}{\thesubsection}{0.6em}{}

\newcolumntype{C}{>{\centering\arraybackslash}X}
\renewcommand{\arraystretch}{1.15}

\setlength{\parskip}{4pt}
\setlength{\parindent}{0pt}

\begin{document}

%---------------------------------------------------------------
\begin{center}
  {\Large\bfseries Adversarial-Robust GNNs for Cyber-Physical\\[2pt]
   Anomaly Detection --- Progress Report}\\[6pt]
  {\normalsize Onur Ozan S\"unger \quad\textbar\quad MSc Thesis, Sapienza
   Universit\`a di Roma}\\[2pt]
  {\small Advisors: Tiziana Cattai, Pierluigi Locatelli \quad\textbar\quad
   \today}
\end{center}

\vspace{2pt}
\hrule
\vspace{6pt}

\noindent
\textbf{Scope.} This report summarises the method and results in two
parts. \textbf{Part~1} develops and validates the detector on
\emph{water distribution networks}. \textbf{Part~2} tests whether the
identical method transfers to two other cyber-physical graphs ---
\emph{power grids} and \emph{road-traffic networks}.

%===============================================================
\section{Part 1 --- Water Distribution Networks}

\subsection{Problem and threat model}
A water network is modelled as a graph $G=(V,E)$: nodes are junctions
with a pressure time-series, edges are pipes with a flow time-series.
Sensors are sparse and noisy (50\% missing, additive noise). An attacker
compromises a subset of sensors and falsifies their readings. Five attack
families are considered, spanning the standard knowledge/capability
spectrum:

\begin{center}\small
\begin{tabularx}{\linewidth}{@{}l l l X@{}}
\toprule
\textbf{Attack} & \textbf{Knowledge} & \textbf{Goal} & \textbf{Signature} \\
\midrule
Random    & black-box & disruption & implausible independent values \\
Noise     & black-box & disruption & inflated high-frequency variance \\
Targeted  & grey-box  & localised damage & bias on a chosen node \\
Stealthy  & white-box & evasion & slow drift within physical envelope \\
Replay    & grey-box  & evasion & re-broadcast of a past legitimate reading \\
\bottomrule
\end{tabularx}
\end{center}

The detector's job, per snapshot, is twofold: \emph{reconstruct} the true
pressure/flow at every node, and \emph{flag} which sensors are
compromised.

\subsection{Architecture}
The model is a \textbf{temporal Mixture-of-Experts Graph Neural Network},
built up in four stages, each removing a specific weakness of the
previous one:

\begin{enumerate}
\item \textbf{Spatial GNN.} A GraphSAGE backbone with four prediction
  heads (pressure/flow reconstruction + anomaly logits) and a
  mass-conservation penalty $\lVert B\mathbf{q}\rVert^2$ ($B$ = incidence
  matrix) that enforces physical consistency.
\item \textbf{Temporal GRU.} A 6-step sliding window with a GRU adds the
  memory needed to reason about \emph{how} a reading evolves --- essential
  for attacks that are plausible in a single snapshot.
\item \textbf{Mixture of Experts.} Six attack-specialised expert networks
  plus a learned \emph{router} (a small GNN) that classifies the dominant
  attack in the window and mixes the experts accordingly. A load-balancing
  entropy term prevents expert collapse.
\item \textbf{Pattern features.} Per-node autocorrelation, adjacent-diff
  standard deviation and noise ratio, which expose the missing observation
  noise that a replayed value lacks.
\end{enumerate}

Four refinements sharpen the routing and expert training:
(i)~\emph{confidence-gated rerouting} --- when the router is unsure the
softmax temperature rises so the prediction is not bet on one possibly
wrong expert; (ii)~a \emph{smaller router with larger experts}
(377\,K parameters); (iii)~\emph{direct expert supervision} --- every
expert is trained on its own attack class regardless of routing, so a
mis-routed class never starves; (iv)~\emph{per-expert reconstruction}.

\newpage
\subsection{Results (Modena network, 272 nodes, 10 seeds)}
Detection is strong and highly reproducible: overall pressure
$\mathbf{F1 = 0.866 \pm 0.004}$, AUROC $0.968$.

\begin{center}\small
\begin{tabularx}{\linewidth}{@{}l CCCCC@{}}
\toprule
 & Random & Targeted & Stealthy & Noise & Replay \\
\midrule
F1 (per attack) & 0.998 & 0.996 & 0.943 & 0.910 & \textbf{0.019} \\
\bottomrule
\end{tabularx}
\end{center}

Two findings are worth highlighting:

\textbf{Router diagnostic.} Measuring the router directly exposes a failure
mode. On Modena it is healthy (96\% accuracy). On the larger Net3 topology
it mis-routes \textbf{77\% of stealthy windows to the replay expert}: both
classes look smooth/low-noise, and 97 nodes are too few to separate them.
The confidence-gated rerouting and direct supervision are designed to
contain exactly this case.

\textbf{The replay ceiling (an honest negative).} Replay resists
detection not for lack of training but for an information-theoretic
reason: over a six-hour window Modena pressure varies by
$\sim$0.01\,m while sensor noise is $\sim$0.4\,m, so a replayed reading is
statistically indistinguishable from the truth. A 10-seed sweep that
weights replay more heavily only trades overall F1 away along a Pareto
front (F1 $0.866\!\to\!0.715$ to lift replay $0.02\!\to\!0.09$). The model
reaches the ceiling; beating it needs a change of \emph{domain} or of
\emph{attacker}, addressed next.

%===============================================================
\section{Part 2 --- Other Cyber-Physical Networks}

\subsection{Method: one detector, three physics}
To test generality the model, the training script and the five-attack
corruption pipeline are kept \textbf{unchanged}; only a data adapter is
added per domain. Each domain maps onto the same
\texttt{(nodes, edges, node-signal, edge-signal)} abstraction:

\begin{center}\small
\begin{tabularx}{\linewidth}{@{}l X X X@{}}
\toprule
 & \textbf{Water (Modena)} & \textbf{Power (IEEE-118)} & \textbf{Traffic (200 sensors)}\\
\midrule
Node / edge   & junction / pipe & bus / line & detector / road \\
Node signal   & pressure (m)    & voltage (pu) & speed (km/h) \\
Edge signal   & flow            & active power & volume \\
Simulator     & EPANET          & pandapower power-flow & rush-hour + diffusion \\
Conservation  & mass            & Kirchhoff (KCL) & flow at junctions \\
\bottomrule
\end{tabularx}
\end{center}

Power data are generated from the IEEE 118-bus system with time-varying
loads; traffic data from a synthetic road-sensor network with realistic
rush-hour speed dynamics and spatial correlation.

\subsection{Results: the detector is domain-agnostic}
Overall detection F1 lands in a tight band across three unrelated physical
systems, with no domain-specific tuning:

\begin{center}\small
\begin{tabularx}{\linewidth}{@{}l C C C C C C@{}}
\toprule
\textbf{Domain} & \textbf{F1} & Random & Targeted & Stealthy & Noise & \textbf{Replay}\\
\midrule
Water (10 seeds)   & $0.866$ & 0.998 & 0.996 & 0.943 & 0.910 & \textbf{0.019} \\
Power (5 seeds)    & $0.845$ & 1.000 & 1.000 & 0.964 & 0.980 & \textbf{0.265} \\
Traffic (5 seeds)  & $0.870$ & 0.987 & 0.992 & 0.658 & 0.754 & \textbf{0.871} \\
\bottomrule
\end{tabularx}
\end{center}

\newpage
\subsection{Key insight: the replay ceiling is a property of the domain}
The most informative result is the \textbf{replay} column. The same model
reports replay F1 of $0.02 \to 0.27 \to 0.87$ across water, power and
traffic. The reason is purely physical: replay is detectable only when the
monitored signal changes appreciably between readings. Water pressure and
grid voltage move slowly, so a stale value hides inside the noise band;
traffic speed swings sharply at rush hour, so a stale value stands out.
This turns the one-domain negative result of Part~1 into a
\emph{general principle}: the replay ceiling is set by the signal
dynamics, not by the detector.

A complementary observation: the \emph{difficulty profile inverts}.
Water and power find stealthy/noise easy but replay impossible; traffic is
the mirror image --- replay is easy, while a stealthy drift blends into
natural congestion. The same detector, opposite hard cases.

%===============================================================
\section{Summary and next steps}

\textbf{Current status.} Part~1 is complete and reproducible
(F1 $0.866\pm0.004$, 10 seeds), with the router failure diagnosed and the
replay ceiling characterised analytically. Part~2 shows the method is
domain-agnostic on power and traffic, and yields a general principle about
replay detectability.

\textbf{Next steps.} (1)~Formalise the cross-domain study as a
\emph{robustness certificate} --- the self-play attacker/defender
equilibrium gives a lower bound on detector performance under adaptive
attack. (2)~Extend the threat-model section and relate it to the
smart-grid FDIA and intelligent-transport-security literature.
(3)~Consolidate the three domains into a single evaluation and identify a
target venue.

\end{document}
